% note éditée à la main (coupe du 17/11/2026) : graphiques.fin ne l'écrase pas
Source : Insee, comptes nationaux 2025 (tableau 8.212) ; DGFiP Statistiques n° 51 ; CPO (2025) ; calculs de l'auteur (fiscalite/is/calculs\_note\_is.py, mécanique de fiscalite/is/simulation\_is.py).\par Note : coût annuel pour l'État, hors intérêts, dans le cas central (75 \% de l'investissement fait par des sociétés imposables, régime actuel pour moitié linéaire, pour moitié dégressif, plafond couvrant 20 \% de l'investissement), sans investissement supplémentaire suscité. Ressources de l'assiette : 3,5 à 4,9 Md\euro{} en 2027, puis +3 \% par an. Les trois calendriers ont le même coût en régime permanent (1,3 à 2,0 Md\euro{} en 2034).
